\documentclass[11pt,a4paper]{article}
\usepackage[utf8]{inputenc}
\usepackage[T1]{fontenc}
\usepackage[brazilian]{babel}
\usepackage{geometry}
\usepackage{graphicx}
\usepackage{booktabs}
\usepackage{float}
\usepackage{enumitem}
\usepackage{amsmath}
\usepackage{xcolor}
\usepackage{tcolorbox}
\geometry{margin=2.5cm}
\setlength{\parindent}{0pt}
\setlength{\parskip}{0.5em}
\tcbuselibrary{skins,breakable}
\newtcolorbox{nota}[1][]{colback=yellow!10,colframe=orange!60!black,fonttitle=\bfseries,title=#1,breakable}

\begin{document}
\begin{center}
  
\includegraphics[width=0.16\textwidth]{../../../figures_global/unb.png}\\[0.3cm]
  {\large\textbf{GABARITO DO PROFESSOR -- USO RESTRITO}}\\[0.2cm]
  {\large FGA0092 -- Princípios de Comunicação para Engenharia}\\[0.2cm]
  \rule{12cm}{0.8pt}\\[0.2cm]
  {\Large\textbf{Prática 06: Demodulação FM}}\\[0.2cm]
  \rule{12cm}{0.8pt}
\end{center}

\section*{Comportamento esperado no Time Sink}

Com $\Delta f = 25\,000$\,Hz, $f_m = 1\,000$\,Hz e \texttt{noise\_amp} $= 0$:

\begin{itemize}
  \item Mensagem original: cosseno limpo de 1\,000\,Hz, amplitude unitária.
  \item Slope detector: cosseno reconhecível porém com distorção harmônica
    visível e amplitude reduzida em relação à mensagem.
  \item Discriminador balanceado: cosseno mais limpo, com menor distorção
    e nível CC quase nulo antes do \texttt{DC Blocker}.
  \item \texttt{WBFM Receive}: cosseno praticamente idêntico à mensagem
    original, com fase e amplitude consistentes.
\end{itemize}

\section*{Verificação via GNU Radio}

A figura abaixo foi gerada rodando o flowgraph equivalente ao
\texttt{pratica\_06\_fm\_parte2.grc} de forma headless via
\texttt{verificar\_gnuradio.py}. Confirma o comportamento esperado e fornece
valores numéricos de referência para correção.

\begin{center}
  \includegraphics[width=0.95\textwidth]{verificacao_gnuradio.png}
\end{center}

\textbf{Resultados medidos} ($f_m = 1$\,kHz, $\Delta f = 25$\,kHz,
$\beta = 25$, sem ruído):

\begin{tabular}{l c c}
  \toprule
  Demodulador & Fidelidade $|\rho|$ & THD \\
  \midrule
  Slope detector simples         & $0{,}7235$ & $94{,}9\%$ \\
  Discriminador balanceado       & $0{,}9701$ & $24{,}5\%$ \\
  \texttt{WBFM Receive}          & $0{,}9999$ & $0{,}3\%$ \\
  \bottomrule
\end{tabular}

\textbf{Robustez ao ruído} ($|\rho|$ alinhada vs \texttt{noise\_amp}):

\begin{tabular}{c c c c}
  \toprule
  \texttt{noise\_amp} & Slope & Balanceado & \texttt{WBFM Receive} \\
  \midrule
  $0{,}00$ & $0{,}724$ & $0{,}970$ & $1{,}000$ \\
  $0{,}05$ & $0{,}718$ & $0{,}970$ & $1{,}000$ \\
  $0{,}10$ & $0{,}712$ & $0{,}970$ & $1{,}000$ \\
  $0{,}20$ & $0{,}700$ & $0{,}969$ & $1{,}000$ \\
  $0{,}50$ & $0{,}670$ & $0{,}964$ & $0{,}999$ \\
  $1{,}00$ & $0{,}627$ & $0{,}926$ & $0{,}895$ \\
  \bottomrule
\end{tabular}

\begin{nota}[Interpretação dos valores medidos]
O slope detector simples tem $|\rho| \approx 0{,}72$ mesmo sem ruído porque
a distorção harmônica intrínseca da inclinação não-linear do filtro é o
fator limitante, não o ruído. O discriminador balanceado e o
\texttt{WBFM Receive} suportam \texttt{noise\_amp} até cerca de $0{,}5$
sem degradação significativa, e degradam de forma visível somente em
\texttt{noise\_amp} $\approx 1{,}0$, onde o \text{SNR}$_i \approx 0$\,dB e o
limiar de FM é cruzado.
\end{nota}

\section*{Respostas esperadas}

\begin{enumerate}[label=\textbf{Q\arabic*.}, leftmargin=*, itemsep=1em]

  \item \textbf{Como o filtro passa-faixa converte variação de frequência em variação de amplitude?}

    \textbf{Resposta esperada:} A resposta em frequência $|H(f)|$ do filtro
    passa-faixa não é constante: na borda da banda de passagem ela varia
    monotonicamente com $f$. Quando a frequência instantânea $f_i(t)$ do
    sinal FM percorre essa região de inclinação, a amplitude da saída do
    filtro acompanha $|H(f_i(t))|$, ou seja, varia proporcionalmente à
    mensagem $m(t)$ que controla $f_i$. O detector de envelope converte essa
    AM em sinal de banda base. Quando $\Delta f$ é pequeno comparado à
    largura do filtro, $f_i(t)$ permanece quase toda na região plana de
    $|H(f)|$, e a conversão FM-AM se torna desprezível, resultando em sinal
    recuperado de amplitude muito baixa.

    \begin{nota}[Erros comuns]
      Afirmar que o filtro ``elimina o ruído'' ou ``recupera diretamente a
      mensagem'' sem explicar o mecanismo de conversão FM-AM por inclinação
      da resposta em frequência.
    \end{nota}

  \item \textbf{Distorção observada no slope detector e sua origem.}

    \textbf{Resposta esperada:} A forma de onda recuperada apresenta
    assimetria entre semiciclos positivos e negativos e harmônicos visíveis,
    porque a curva de transferência $|H(f)|$ não é perfeitamente linear ao
    longo da faixa $[f_c - \Delta f, f_c + \Delta f]$. Termos quadráticos e
    de ordem superior do polinômio que aproxima $|H(f)|$ geram distorção
    harmônica em $2 f_m$, $3 f_m$ etc. Além disso, há um nível CC residual
    proporcional ao valor médio de $|H(f_i(t))|$, que precisa ser removido
    pelo \texttt{DC Blocker}.

    \begin{nota}[Erros comuns]
      Atribuir a distorção apenas ao ruído ou à decimação, sem identificar
      a não-linearidade da resposta do filtro como causa principal.
    \end{nota}

  \item \textbf{Vantagens do discriminador balanceado.}

    \textbf{Resposta esperada:} Os dois ramos produzem envoltórias
    aproximadamente $K_0 + K_1 m(t) + \alpha m^2(t) + \cdots$ e
    $K_0 - K_1 m(t) + \alpha m^2(t) + \cdots$. A subtração cancela o termo
    constante $K_0$ e os termos de ordem par, restando $2 K_1 m(t)$ mais
    distorção de ordem ímpar, conforme a Equação~(2) do enunciado. O
    resultado é maior fidelidade na recuperação da mensagem, melhor
    linearidade na região central e dispensa parcial do bloqueio CC. No
    Time Sink isso aparece como um cosseno mais simétrico, com menos
    distorção visível e amplitude maior.

    \begin{nota}[Erros comuns]
      Dizer apenas que ``dois filtros são melhores que um'' sem mostrar
      explicitamente o cancelamento do termo CC e da distorção de ordem par.
    \end{nota}

  \item \textbf{Comparação dos três demoduladores e por que \texttt{WBFM Receive} costuma vencer.}

    \textbf{Resposta esperada:} A ordem esperada de fidelidade é
    \texttt{WBFM Receive} $>$ discriminador balanceado $>$ slope detector.
    O \texttt{WBFM Receive} usa um discriminador de quadratura, que estima
    $f_i(t)$ a partir da fase instantânea do sinal complexo, sendo
    teoricamente linear em frequência ao longo de toda a faixa de operação.
    Os detectores baseados em envoltória dependem de aproximar uma região
    linear na resposta de filtros reais, o que introduz distorção residual.
    Adicionalmente, o discriminador de quadratura tem ganho independente da
    forma de onda do filtro, tornando-o mais robusto a tolerâncias de
    componentes.

    \begin{nota}[Erros comuns]
      Afirmar que o \texttt{WBFM Receive} é melhor por ``ser do GNU Radio''
      sem descrever o princípio de operação do discriminador de quadratura.
    \end{nota}

  \item \textbf{Limiares de \texttt{noise\_amp} e efeito de limiar do FM.}

    \textbf{Resposta esperada:} A simulação mostra que o fator limitante do
    slope detector é a distorção intrínseca, não o ruído: $|\rho| \approx 0{,}72$
    já em \texttt{noise\_amp} $= 0$ e cai pouco com ruído adicional. O
    discriminador balanceado e o \texttt{WBFM Receive} suportam
    \texttt{noise\_amp} até cerca de $0{,}5$ ($\text{SNR}_i \approx 6$\,dB) sem
    degradação significativa em $|\rho|$, e só apresentam queda visível em
    \texttt{noise\_amp} $\approx 1{,}0$ ($\text{SNR}_i \approx 0$\,dB), onde o
    limiar do FM é cruzado e cliques aparecem na saída. A ordenação esperada
    em fidelidade é \texttt{WBFM Receive} $>$ Balanceado $\gg$ Slope, o que
    confirma que demoduladores baseados em fase instantânea superam os
    baseados em envoltória, e que o discriminador balanceado lineariza
    parcialmente a resposta da inclinação.

    \begin{nota}[Erros comuns]
      Reportar valores numéricos sem ordenação relativa entre os três
      demoduladores, ou não mencionar o efeito de limiar do FM.
    \end{nota}

\end{enumerate}

\section*{Critérios de correção}
\begin{tabular}{ll}
  \toprule
  Critério & Pontuação \\
  \midrule
  Q1 -- Mecanismo de conversão FM-AM por inclinação                              & 1,0 ponto \\
  Q2 -- Identificação da distorção harmônica e nível CC residual                 & 1,0 ponto \\
  Q3 -- Cancelamento de CC e distorção par no discriminador balanceado           & 1,0 ponto \\
  Q4 -- Comparação dos três demoduladores com justificativa de princípio         & 1,0 ponto \\
  Q5 -- Ordenação dos limiares de ruído e menção ao limiar de FM                 & 1,0 ponto \\
  \midrule
  \textbf{Total} & \textbf{5,0 pontos} \\
  \bottomrule
\end{tabular}

\bigskip
\begin{nota}[Nota de correção geral]
  Priorizar a coerência física e a evidência experimental por meio de
  capturas de tela do Time Sink. Respostas que confundam detector de
  envelope com discriminador de quadratura devem ter desconto. Valores
  numéricos de limiar de \texttt{noise\_amp} podem variar com a
  implementação dos filtros, e o critério de correção deve focar na
  ordenação relativa entre os demoduladores, não nos valores absolutos.
\end{nota}

\end{document}
